\chapter{Writing strategies}
\label{ch:strategies}

A \textbf{strategy} is whatever Scheme decides how a turn goes. The default
one is a stack of small files; yours can replace it, or, much more often,
wrap one piece of it. This chapter teaches the wrapping idiom, the one rule
that keeps wrappers correct, and a set of strategies you can drop into an
\code{init.scm} today.

\section{How the default loop is layered}

The session loads, in this order:

\begin{enumerate}
  \item the prelude (\code{assoc-ref}, \code{compose}\ldots);
  \item \file{strategy.scm} (\code{nodwyn-turn}, \code{agent-turn},
    \code{ask-model}, the evidence loop), \file{prompts.scm},
    \file{compact.scm}, \file{agents.scm} (\code{run-calls},
    \code{sub-agent}), \file{optimize.scm}, \file{plan.scm};
  \item the working habits, each wrapping what came before:
    \file{shell.scm}, \file{hygiene.scm}, \file{readdef.scm},
    \file{envhint.scm}, \file{checkloop.scm};
  \item your \file{.nodwyn/init.scm}.
\end{enumerate}

Because Scheme looks up a global name when it is \emph{called}, not when
the calling function was defined, redefining \code{run-calls} in step 4
changes what the default \code{agent-turn} from step 2 calls. That is the
whole trick.

\section{The wrapping idiom}

\begin{scm}
(define my-base-run-calls run-calls)       ; 1. keep what is there now
(define (run-calls calls)                  ; 2. define the new one
  ;; ... before ...
  (let ((results (my-base-run-calls calls)))   ; 3. call the old one
    ;; ... after ...
    results))
\end{scm}

Step 1 saves the \emph{current} \code{run-calls}, which may itself be a
habit that wraps another habit that wraps the original. Each layer adds one
behaviour and passes the rest down. \nodwyn's own habits are written
exactly like this, so yours compose with them.

\begin{tip}
Give the saved name a prefix that is yours (\code{my-base-},
\code{tally-base-}). Two layers that both call their saved copy
\code{base-run-calls} would overwrite each other's, and the second layer
would call itself forever.
\end{tip}

\section{The one rule: one result per call, in order}

\code{run-calls} must return exactly one result for each call it was
given, in the same order. The default \code{agent-turn}, the evidence loop
and the habits all walk calls and results side by side; a wrapper that
drops a result shifts every later one onto the wrong call.

So when your wrapper decides \emph{not} to run a call, skip it with a
reason \emph{and} return a result for it. A helper makes that painless:

\begin{scm}
;; Skip one call with a reason, and return the result run-calls would.
(define (skipped call reason)
  (skip-tools (list call) reason)
  (list (cons 'id (assoc-ref call 'id))
        (cons 'name (assoc-ref call 'name))
        (cons 'ok #f)
        (cons 'value (string-append "Not executed: " reason))))
\end{scm}

\section{Strategy: stop a batch at the first failure}

Models often send several calls at once, each assuming the previous one
worked. This runs them one at a time and skips the rest of the batch after
a failure, so the model re-plans from the error:

\begin{scm}
(define stop-base-run-calls run-calls)

(define (run-calls calls)
  (let loop ((cs calls) (out '()) (failed #f))
    (cond ((null? cs) (reverse out))
          (failed
           (loop (cdr cs)
                 (cons (skipped (car cs) "an earlier call in this batch failed") out)
                 #t))
          (else
           (let ((r (car (stop-base-run-calls (list (car cs))))))
             (loop (cdr cs) (cons r out) (not (assoc-ref r 'ok #f))))))))
\end{scm}

It runs the calls one at a time through the saved \code{stop-base-run-calls}
(calling \code{run-calls} there would call itself forever), and after the
first failure answers every remaining call with a skip and a result. In a real session, with
the first call denied:

\begin{txt}
▸ fs-list {"path":"."}
⚑ guard deny fs-list: user checker no-listing
◂ fs-list error: fs-list: denied: user checker no-listing
◂ fs-read error: Not executed: an earlier call in this batch failed
\end{txt}

\section{Strategy: count what each turn used}

\begin{scm}
(define tally '())

(define (tally-add! name)
  (let ((seen (assoc name tally)))
    (set! tally
          (if seen
              (map (lambda (p) (if (equal? (car p) name) (cons name (+ 1 (cdr p))) p))
                   tally)
              (cons (cons name 1) tally)))))

(define tally-base-run-calls run-calls)
(define (run-calls calls)
  (for-each (lambda (c) (tally-add! (assoc-ref c 'name))) calls)
  (tally-base-run-calls calls))

(define tally-base-agent-turn agent-turn)
(define (agent-turn request)
  (set! tally '())
  (tally-base-agent-turn request)
  (log-note (string-append "tools this turn: "
                           (string-join (map (lambda (p) (format "~a x~a" (car p) (cdr p)))
                                             (reverse tally))
                                        ", "))))
\end{scm}

\begin{txt}
◂ fs-read ok (45 bytes, 6 ms)
...
note: tools this turn: fs-list x1, fs-read x1
\end{txt}

\code{log-note} writes to the log and the transcript; the model never sees
it. It is the right way for a strategy to leave a trail for you.

\section{Strategy: a look-only switch}

A variable you flip with \code{:eval} mid-session, that makes the model
read and report but change nothing:

\begin{scm}
(define look-only #f)
(define look-only-writers '("fs-write" "fs-edit" "shell-exec"))
(define look-only-reason "look-only mode: read and report, change nothing")

(define look-base-run-calls run-calls)
(define (run-calls calls)
  (if (not look-only)
      (look-base-run-calls calls)
      (map (lambda (c)
             (if (member (assoc-ref c 'name) look-only-writers)
                 (skipped c look-only-reason)
                 (car (look-base-run-calls (list c)))))
           calls)))
\end{scm}

\begin{txt}
you> :eval (set! look-only #t)
you> tidy up the notes folder
\end{txt}

The model's writes come back as \code{Not executed: look-only mode\ldots},
which tells it plainly what happened, so it adapts instead of retrying.

\begin{careful}
Switches like this are conveniences, not security. The guard is the
security: a strategy runs with harness origin and can only \emph{skip}
calls. For a rule that must hold, write a \code{register-check} with
\code{"deny"}.
\end{careful}

\section{Strategy: don't stop until the tests ran}

The default loop already does this when nobody is watching (the
\emph{evidence loop}, Chapter~\ref{ch:plans}). Here is a small version for
an interactive session, to see how \code{add-message} works:

\begin{scm}
(define (ran-tests? calls)
  (let any ((cs calls))
    (and (pair? cs)
         (or (and (equal? (assoc-ref (car cs) 'name) "shell-exec")
                  (string-contains? (assoc-ref (car cs) 'args "") "test"))
             (any (cdr cs))))))

(define (agent-turn request)
  (let loop ((reply (ask-model)) (tested #f) (nudged #f))
    (let ((calls (assoc-ref reply 'tool-calls '())))
      (cond ((pair? calls)
             (run-calls calls)
             (maybe-compact request reply)
             (loop (ask-model) (or tested (ran-tests? calls)) nudged))
            ((and (not tested) (not nudged))
             (add-message "Before you finish, run the tests that cover your change and report the result.")
             (loop (ask-model) tested #t))))))
\end{scm}

\code{add-message} speaks to the model as a user message and is logged as
\code{harness.msg}. It is refused while calls from the latest reply are
still unanswered, because a call must be followed by its result.

This \code{agent-turn} replaces the default one, so it also drops the plan
gate of an interactive turn (a run's turns still plan, since
\code{nodwyn-turn} runs the gate first) and the \code{/compact} and
\code{/optimize} handling. That is fine for an experiment. For keeps, wrap
instead:

\begin{scm}
(define mine-base-agent-turn agent-turn)
(define (agent-turn request)
  (if (string-prefix? "/review" request)
      (review-turn request)              ; your own command
      (mine-base-agent-turn request)))
\end{scm}

\section{\texttt{ask-model} and cut-off replies}

A reply that reaches the output limit comes back with \code{done} set to
\code{"length"}. If it had no tool call to answer, the turn would simply
end with half an answer. \code{ask-model} is \code{model-turn} with a
remedy: it adds a note asking the model to be briefer and asks again, at
most \code{output-cut-retries} (2) times. Use \code{ask-model} wherever you
would use \code{model-turn}; the default loop, the planner and sub-agents
all do.

\begin{puzzle}[ 15: read before write]
Write a \code{run-calls} wrapper that skips an \code{fs-edit} of a file the
model has not read with \code{fs-read} earlier in the session, with the
reason \code{"read the file before editing it"}. Keep the set of read paths
in a variable. Remember the one rule.\ans{15}
\end{puzzle}
